\documentclass[11pt]{article}

\input{../../preamble.tex}
% --- Page geometry tuned to match the PDF look ---
\usepackage[margin=1in]{geometry}
\usepackage{amsmath,amssymb,amsfonts}
\usepackage{enumitem}
\usepackage{mathtools}

\setlength{\parindent}{0pt}
\setlength{\parskip}{6pt}

\newcommand{\ellinfty}{\ell^{\infty}}
\newcommand{\ellone}{\ell^{1}}
\newcommand{\ellzero}{\ell_{0}}
\newcommand{\ellc}{\ell_{c}}
\newcommand{\norm}[1]{\left\lVert #1 \right\rVert}
\newcommand{\restr}[2]{#1\!\upharpoonright_{#2}}


\begin{document}

\begin{center}
{\LARGE Math 421/510: Homework 2}
\end{center}

\vspace{1.6em}

\begin{enumerate}
\item Let $\ell_c$ denote the space of finite sequences, endowed with the supremum norm $\|\cdot\|_\infty$.
Define $T:\ell_c \to \ell_c$ by
\[
T:(x_1,x_2,x_3,\ldots)\mapsto \left(x_1,\frac{1}{2}x_2,\frac{1}{3}x_3,\ldots\right).
\]
Show that $T$ is well-defined, linear, bounded, and bijective, but that $T^{-1}$ is not bounded. Why does this not contradict the open mapping theorem?


\textbf{Solution.}

First, we confirm $T$ is well defined. 
For $x \in \ell_{c}$, $Tx$ is a well defined sequence by definition, so it remains to check if $Tx \in \ell_c$.
Since $x$ is a finite sequence, there exists $N$ such that $x_n = 0$ for $n > N$.
However, since $(Tx)_n = \frac{1}{n} x_n$, $x_n = 0 \implies (Tx)_n = 0$ so then $Tx$ is also a finite sequence, so $Tx \in \ell_c$.

For linearity, we see that
\[
T(x + \lambda y) = (x_1 + \lambda y_1, \frac{1}{2} x_2 + \frac{\lambda}{2} y_2, \ldots) = (x_1, \frac{1}{2} x_2, \ldots) + \lambda (y_1, \frac{1}{2} y_2, \ldots) = Tx + \lambda Ty
.\] 
Therefore, $T$ is linear.

To show boundedness, consider $\|x\|_{\infty} = 1$.
Then, $|(Tx)_n| = \frac{1}{n} |x_n| \leq \frac{1}{n} \leq 1$.
Thus, $\|Tx\|_{\infty} = \sup_{n}|(Tx)_n| \leq 1$, and since this applies for all $x$ with $\|x\|_{\infty} = 1$ it follows that $\|T\| \leq 1$.


To show $T$ is a bijection, we show injectivity and surjectivity.
Suppose $Tx = 0$.
Then, $\frac{1}{n}x_n = 0 \implies x_n = 0 \forall n \in \N \implies x = 0$.
Now, suppose $\{y_n\} \in \ell_{c}$.
Then, if we define $\{x_n\} $ such that $x_n = ny_n$, we have $(Tx)_{n} = \frac{1}{n} n y_n = y_n$, so $Tx=y$.
This also gives us that $T^{-1}(x_1, x_2, x_3, \ldots) = (x_1, 2x_2, 3x_3, \ldots)$.

Finally, we show that $T^{-1}$ is not bounded.
Consider a sequence of $n$ 1s, that is, $x = (1, \ldots, 1, 0, \ldots)$.
Then, $T^{-1}(x) = (1, 2, 3, \ldots, n, 0 ,0, \ldots)$.
But then, $\|T^{-1}(x)\|_{\infty} = n$.
Therefore, for all $n$, there exists some $x \in \ell_{c}$ with $\|x\|_{\infty} = 1$ such that $\|T^{-1}x\|_{\infty} = n$, so $T^{-1}$ is unbounded.

This does not contradict the Corollary to the open mapping theorem, because $\ell_{c}$ is not a banach space (that is, it is not complete) w.r.t. the supremum norm.
To see this, consider that the sequence $y^{(n)} = (1, \ldots, \frac{1}{n}, 0, \ldots)$ is Cauchy (indeed, if w.l.o.g. $m > n$, then $\|y^{(n)} - y^{(m)}\| = \sup_{i} |y^{(n)}_i - y^{(m)}_i| = 1 / (n+1)$), but its limit is an infinite sequence, so $\lim_{n \to \infty}y^{(n)} \not\in \ell_{c}$.
Therefore, the requirements of the Open Mapping Theorem that $\mathcal{X}$ and $\mathcal{Y}$ are Banach spaces is not fulfilled, so we can't expect the theorem to hold.

\item Let $Y\subset \ell^1$ denote the space of sequences $x=(x_1,x_2,\ldots)$ for which
$\sum_{n=1}^\infty n|x_n|<\infty$, endowed with the norm
$\|x\|_1=\sum_{n=1}^\infty |x_n|$.
Define the linear map $T:Y\to \ell^1$ by
\[
(Tx)_n = n x_n.
\]
Show that $T$ is closed, but not bounded. Why does this not contradict the closed graph theorem?


\textbf{Solution.}

We show that $T$ is closed.
Suppose that $x^{(n)} \in Y$ is a sequence such that $x^{(n)} \to x \in Y$, and suppose that we know that $Tx^{(n)} \to y \in \ell^{1}$.
For $T$ to be closed, we need to show that $Tx = y$.
Fix $m \in \N$.
Then, we have that
\begin{align*}
    |(Tx)_m - (y)_m| &\leq |(Tx)_m - (Tx^{(n)})_m| + |(Tx^{(n)})_m - y_m|\\
                     &\leq m| (x)_m - (x^{(n)})_m | + \|Tx^{(n)} - y\|_{1}\\
                     &\leq m \|x - x^{(n)}\|_1 + \|Tx^{(n)} - y\|_1\\
                     &\to 0 \text{ as } n \to \infty
.\end{align*}
Therefore, $(Tx)_m = y_m$ for all $m$, so $Tx = y$. Thus, $T$ is closed.

However, $T$ is not bounded.
Consider $x^{(n)} = (0, \ldots,0, 1, 0, \ldots)$, where the $1$ is in the $n$-th location.
Then, $\|x^{(n)}\|_1 = 1$, and $\sum_{m=1}^{\infty}m |(x^{(n)})_m| = n < \infty$, so $x^{(n)} \in Y$.
However, $\|T x^{(n)}\|_1 = n$.
Therefore, $\|T\| \geq n$ for all $n \in \N$, so we have that $T$ is not bounded.

This does not contradict the closed graph theorem, because $Y$ is not a Banach space. To show, this, we construct a sequence in $Y$ which is cauchy that converges to an element in $\ell^{1} \setminus Y$.
Consider the sequence $x_n = \frac{1}{n^2}$, and let $x^{(N)}$ be the sequence such that $x^{(N)}_n = 0$ if $n > N$, and $x^{(N)}_n = (x)_n$ if $n \le N$.
Then, since $x^{(N)}$ is a finite sequence, $x^{(N)} \in Y$.
However, $\lim_{N \to \infty} x^{(N)} = x$ under the $\|\cdot\|_1$ norm, since for $N > M$, we have
\[
\|x - x^{(N)}\| = \sum_{n = N}^{\infty} \frac{1}{n^2} < \sum_{n=M}^{\infty} \frac{1}{n^2} \to 0 \text{ as } M \to \infty
.\] 

Also, $\|x\| = \sum_{n=1}^{\infty} \frac{1}{n^2} = \pi^2 / 6$, so $x \in \ell^{1}$.
However, $\sum_{n=1}^{\infty}n |(x)_n| = \sum_{n=1}^{\infty} \frac{1}{n} = \infty$, so $x \not\in Y$.
Therefore, we have a sequence in $Y$, $\{x^{(N)}\}_{N=1}^{\infty} $ which converges to $x$, but $x \not\in  Y$.
Therefore, $Y$ is not closed, so it is not complete, so it is not Banach.
Even if $T : Y \to \ell^{1}$ is closed, the hypothesis of the closed graph theorem that both $Y$ and $\ell^{1}$ are banach is not satisfied.

\item ([Folland, ex.\ 5.39]): Let $X$, $Y$ and $Z$ be Banach spaces, and $B:X\times Y\to Z$ be such
that $\forall\, x\in X$, $B(x,\cdot):Y\to Z \in L(Y,Z)$ and $\forall\, y\in Y$, $B(\cdot,y):X\to Z \in L(X,Z)$.
\begin{enumerate}
\item show $\|B(x,y)\|\le C\|x\|\|y\|$ for some $C$;
\item show that $B$ is continuous (as a map from Banach space $X\times Y$ to $Z$.
\end{enumerate}

\textbf{Solution.}

\begin{enumerate}
    \item Fix $x \in X$, with $x \neq 0$.
        Consider the linear map $\frac{B(x, \cdot)}{\|x\|} : \mathcal{Y} \to \mathcal{Z}$.
        Then, since $B(x, \cdot) \in L(Y, Z)$, there exists some constant $C_x$ such that $\|B(x, \cdot)(y)\| \leq C_x \|y\|$, and so $\|\frac{B(x, \cdot)(y)}{\|x\|}\| \leq \frac{C_x}{\|x\|} \|y\|$.

        Now, fix $y \in Y$.
        Since $B(\cdot, y) \in L(X, Z)$, we have that $\|B(x, y)\| \leq \|B(\cdot, y)\| \|x\|$.
        However, this implies that $\frac{B(x, y)}{\|x\|} \leq \|B(\cdot, y)\|$ for all $x \in X$ with $x \neq 0$, so
        \[
            \sup_{x \in X \setminus \{0\} } \|B(x, \cdot)(y) / \|x\|\| \leq \|B(\cdot, y)\|
        .\] 
        
        Since $Y$  is a Banach space, and this holds for all $y \in Y$, the requirements of the Uniform Boundedness Principle are satisfied, so we have that $\sup_{x \in X \setminus \{0\} } \|\frac{B(x, \cdot)}{\|x\|}\| \leq M$ for some $M \in \R$.
        Then, $\|B(x, \cdot)\| \leq M \|x\|$, which by definition means that $\|B(x, y)\| \leq M \|x\|\|y\|$ when $\|x\| \neq 0$.
        When, $x = 0$, then clearly $\|B(0, \cdot)\| = 0 \leq M$.
    \item Let $(x, y) \in X \times Y$.
        Then, from part (a), we have some $M$ such that $\|B(x, y)\| \leq M \|x\| \|y\|$.
        Let $\epsilon > 0$.
        Choose $\delta > 0$ such that $\delta < \min \{\frac{\epsilon}{3M \|x\|}, \frac{\epsilon}{3 M \|y\|}, \frac{\sqrt{\epsilon}}{\sqrt{3M}}\} $.
        Then, when $\|(\Delta x, \Delta y)\| < \delta$ so that $\|\Delta x\|,\| \Delta y\| < \delta$, we have
        \begin{align*}
            \| B(x + \Delta x, y + \Delta y) - B(x, y) \| &=  \|B(x, y) + B(\Delta x, y) + B(x, \Delta y) + B(\Delta x, \Delta y) - B(x, y)\| \\
            &=\leq \|B(\Delta x, y)\| + \|B(x, \Delta y)\| + \| B (\Delta x, \Delta y)\| \\
            &\leq M \|\Delta x\| \|y\| + M \|x\| \|\Delta y\| + M \|\Delta x\| \|\Delta y\|  \\
            &< M \delta \|y\| + M \|x\| \delta + M \delta^2 \\
            &< \frac{\epsilon}{3} + \frac{\epsilon}{3} + \frac{\epsilon}{3} = \epsilon
        .\end{align*}
        Therefore, $B(x, y)$ is continuous as a bilinear map from $X \times  Y$ to $Z$.
\end{enumerate}


\item Suppose $\{g_n\}_{n=1}^\infty \subset L^\infty$ with $\|g_n\|_\infty \le C$, and $g_n\to g$ a.e.
Show $g\in L^\infty$ with $\|g\|_\infty \le C$.


\textbf{Solution.}

We need to show that $g: X \to \R$ is measurable, and $\|g\|_{\infty} < \infty$.
By proposition 2.11 b. in Folland, since $g_n$ is measurable for $n \in \N$, and $g_n \to g$ a.e., $g$ is measurable. 
Now, for $D > C$, we have that
\begin{align*}
    \{x: |g(x)| > D\}  &\subset  \{x:g(x) \neq \lim_{n \to \infty}g_n(x)\} \cup \{x:\lim \sup |g_n(x)| > D\} \\
                    & = \{x:g(x) \neq \lim_{n \to \infty}g_n(x)\} \cup \bigcap_{m = 1}^{\infty} \bigcup_{n = m}^{\infty} \{x: |g_n(x)| > D\} \\
                    &\subset \{x:g(x) \neq \lim_{n \to \infty}g_n(x)\} \cup \bigcup_{n = 1}^{\infty} \{x: |g_n(x)| > D\}
.\end{align*}
But then,
\[
\mu (\{x: |g(x)| > D\} ) \leq \mu(\{x:g(x) \neq \lim_{n \to \infty}g_n(x)\}) +\sum_{n=1}^{\infty} \mu(\{x : |g_n(x)| > D\} ) = 0 + \sum_{n=1}^{\infty} 0
.\] 
Since this applies for all $D > C$, we have that
\[
\inf \{a \geq 0 : \mu(\{x : |g(x)| > a\} ) = 0\}  \leq C
,\] 
as desired. In particular, $\|g\|_{\infty} \leq C$, and $\|g\|_{\infty} < \infty$, so $g \in L^{\infty}$.

\item This exercise gives an another example of a situation where the linear isometry
$L^1 \to (L^\infty)^*$ is not surjective. In assignment \#1, you showed the existence (by invoking
the H-B theorem) of $T\in (\ell^\infty,\|\cdot\|_\infty)^*$, with $\|T\|=1$, and $TS=S$ where
$S:(x_1,x_2,x_3,\ldots)\mapsto (x_2,x_3,\ldots)$ is the shift operator.
Show that $T$ is not given by $Tx=y\cdot x$ for any $y\in \ell^1$.

\textbf{Solution.}

We showed last homework in problem 6 (e) that if $x = (x_1, x_2, \ldots) \in \ell^{\infty}$ satisfies $\lim_{k \to \infty} x_k = L$, then $Tx = L$.
Let $y \in \ell^{1}$ be arbitrary. In particular, $\sum_{n =1 }^{\infty} |y_n| < \infty$.
Then, there exists some $N$ such that $\sum_{n = N}^{\infty} |y_n| < 1$.
Now, consider $x = (x_1, x_2, \ldots)$ such that $x_i = 0$ for $i < N$, and $x_i = 1$ for $i \ge  N$.
In particular, $\lim_{i \to \infty} x_i = 1$, so by 6 (e) from last homework, we have $Tx = 1$.
However, 
\[
|y \cdot x| = \left| \sum_{n = 1}^{\infty} x_n y_n \right| = \left|\sum_{n = N}^{\infty} y_n \right| \leq \sum_{n = N}^{\infty} |y_n| < 1
,\] 
so $y\cdot x \neq 1 = Tx$.
Since $y \in \ell^{1}$ was arbitrary, $T$ is not given by $y \cdot x$ for any $y \in \ell^{1}$.

\end{enumerate}

\end{document}

